\section{Lean}\label{sec:nlean}

Les trois arches du \S\ref{sec:morearches} sont démontrées dans le paquet Lake \texttt{lean/bridges} du dépôt de cette édition (\url{https://github.com/Commutator-IO/grothendieck-project}), qui dépend de lean-pool à une révision fixée; c'est là aussi que sont formalisés les corollaires sur les tabliers des notes de Grothendieck. Les cinq arches de \lp{}, dans lean-pool, sont démontrées en Lean~4 avec Mathlib, sans autre axiome que \texttt{propext}, \texttt{Classical.choice} et \texttt{Quot.sound}. Cette affirmation décrit le développement: les fichiers ne contiennent aucune commande \texttt{\#print axioms}, et nous ne consignons aucune vérification séparée à leur sujet. Pour \texttt{lean/bridges}, l'intégration continue du dépôt de cette édition fait un audit des axiomes à chaque modification: à la version décrite ici, il portait sur 437 déclarations, toutes limitées à \texttt{propext}, \texttt{Classical.choice} et \texttt{Quot.sound}. Les fichiers du paquet utilisés ici sont les suivants; les autres servent à l'article associé \cite{bridges-generator}.
\begin{itemize}
\item \lean{PreorderPresheaves.lean} (\lean{exists_preorder_moritaEquivalentCat_iff_isThin}, d'après l'exercice de Grothendieck, dossier 161-3, p.~26); \lean{SubterminalArches.lean} (les arches deux valeurs et connexe: \lean{forall_subfunctor_terminal_eq_bot_or_eq_top_iff}, \lean{isConnected_iff_forall_isCompl}); \lean{GrothendieckDecks.lean} (les tabliers des notes de Grothendieck: \lean{simplexCategory_twoValued}, \lean{simplexCategory_isConnected}, \lean{simplexCategory_not_spatial}, \lean{singleObj_twoValued}, \lean{singleObj_isConnected}, \lean{singleObj_spatial_iff}, \lean{preorder_twoValued_iff}).
\end{itemize}
Ne sont pas formalisés l'équivalence entre le topos de préfaisceaux d'un préordre et le topos d'un espace (pont~\ref{br:spatial}), le second exercice du dossier 161-3, p.~26 (\S\ref{sec:precedents}), et les énoncés du \S\ref{sec:pseudofields}.

\begin{center}\scriptsize\emergencystretch=3em
\begingroup\def\UrlBreaks{\do\_\do\.\do\/}
\begin{tabularx}{\linewidth}{>{\raggedright\arraybackslash}p{2.6cm}>{\raggedright\arraybackslash}X>{\raggedright\arraybackslash}p{2.9cm}>{\raggedright\arraybackslash}p{2.2cm}}
\toprule
énoncé & déclaration Lean & fichier & statut\\
\midrule
petites catégories: topos de préfaisceaux d'un préordre si et seulement si mince (Grothendieck, 161-3, p.~26) & \lean{exists_preorder_moritaEquivalentCat_iff_isThin} & \lean{PreorderPresheaves} & démontré en Lean\\
arches deux valeurs et connexe (ponts~\ref{br:twoval} et~\ref{br:conn}) & \lean{forall_subfunctor_terminal_eq_bot_or_eq_top_iff}, \lean{isConnected_iff_forall_isCompl} & \lean{SubterminalArches} & démontré en Lean\\
tabliers des notes de Grothendieck: $\Delta$, groupes, préordres (\S\ref{sec:grdecks}) & \lean{simplexCategory_twoValued}, \lean{simplexCategory_isConnected}, \lean{simplexCategory_not_spatial}, \lean{singleObj_twoValued}, \lean{singleObj_isConnected}, \lean{singleObj_spatial_iff}, \lean{preorder_twoValued_iff} & \lean{GrothendieckDecks} & démontré en Lean\\
\bottomrule
\end{tabularx}
\endgroup
\end{center}
